\subsection{埃尔弗--舍斯特兰德公式}

在本节中，E 表示一个复希尔伯特空间。我们将对自伴部分算子的函数演算的某些情形得到一个公式，它直接用所论算子的预解式来表达。

我们给 R （相应地，C）赋以勒贝格测度，记作 \(\mu\) ，并用映射 \((x,y)\mapsto\exp(2i\pi xy)\) 把群 R 与它的对偶群等同起来（参见 II, p. 236 的推论 3）。

对函数 \(f\) ——它在 \(\mathbf{R}^{2}\) 的开集 U 上定义且可微，该开集等同于 \(\mathbf{C}\) ，赋以实坐标 \(x\) 与 \(y\) ——我们记

\[
{\frac {\partial f}{\partial \overline {{{z}}}}} = {\frac {1}{2}} {\Bigl (} {\frac {\partial f}{\partial x}} + i {\frac {\partial f}{\partial y}} {\Bigr)}
\]

（参见 VAR, R2, 8.8.10, p. 24）。

引理 5. --- 设 \(f \in \mathcal{D}(\mathbf{R})\) 。存在函数 \(\tilde{f}\) 属于 \(\mathcal{D}(\mathbf{C})\) ，它与 \(f\) 在 \(\mathbf{R}\) 上重合，且满足

\[
\frac {\partial \widetilde {f}}{\partial \overline {{z}}} (x, 0) = 0\tag{11}
\]

对所有 \(x \in \mathbf{R}\) 成立。这时我们有

\[
\frac {\partial \widetilde {f}}{\partial y} (x, 0) = i f ^ {\prime} (x)\tag{12}
\]

对所有 \(x \in R\) 成立，且存在实数 C ≥ 0 使得

\[
\left| \frac {\partial \widetilde {f}}{\partial \overline {{z}}} (x, y) \right| \leqslant C | y |\tag{13}
\]

对所有 \((x,y)\in \mathbf{R}^2\)

存在 \(\varphi\in\mathcal{D}(\mathbf{R})\) ，其支撑包含在 \([-2,2]\) 中，且在 \([-1,1]\) 上等于 1（IV, p. 196 的引理 1）。置

\[
\widetilde {f} (x, y) = \bigl (f (x) + i y f ^ {\prime} (x) \bigr) \varphi (y)
\]

对 \((x,y)\in\mathbf{R}^{2}\) 。我们有 \(\widetilde{f}\in\mathcal{D}(\mathbf{C})\) ，且 \(\widetilde{f}\) 在 R 上与 f 重合。此外，对任意 \((x,y)\in\mathbf{R}^{2}\) ，可得

\[
\frac {\partial \widetilde {f}}{\partial \overline {{z}}} (x, y) = \frac {1}{2} \Big ((i f (x) - y f ^ {\prime} (x)) \varphi^ {\prime} (y) + i y f ^ {\prime \prime} (x) \varphi (y) \Big).
\]

由于函数 \(\varphi\) 在 0 的一个邻域中等于 1，我们有 \(\varphi'(0) = 0\) ，由此得 (11)。

设 \(\tilde{f}\) 是 \(\mathcal{D}(\mathbf{C})\) 中满足 (11) 的函数。由此即得公式 (12)，而估计式 (13) 则借助中值定理（FVR, I, p. 23, 定理 2）得到。

我们称 \(\tilde{f}\) 是 f 的一个几乎解析延拓。

引理 6. --- 设 \(\varepsilon > 0\) 。函数 \(\sigma_{\varepsilon}\) 定义在 \(\mathbf{R}\) 上，由下式给出

\[
\sigma_ {\varepsilon} (x) = \frac {2 i \varepsilon x}{x ^ {2} + \varepsilon^ {2}}
\]

属于 \(L^{2}(\mathbf{R})\) 。它的傅里叶变换是 \(L^{2}(\mathbf{R})\) 中函数 \(\eta_{\varepsilon}\) 的类，该函数在 0 处取值 0 且满足

\[
\eta_ {\varepsilon} (y) = \frac {2 \pi \varepsilon y}{| y |} e ^ {- 2 \pi \varepsilon | y |}
\]

对所有 \(y \neq 0\) 成立。

由 IV, p. 199 的命题 3，有 \(\sigma_{\varepsilon} \in \mathrm{L}^{2}(\mathbf{R})\) ，且函数 \(\eta_{\varepsilon}\) 的类属于 \(\mathrm{L}^{2}(\mathbf{R}) \cap \mathrm{L}^{1}(\mathbf{R})\) 。对每个 \(x \in \mathbf{R}\) ，我们有

\[
\begin{array}{r} \overline {{\mathcal {F}}} (\eta_ {\varepsilon}) (x) = 2 \pi \varepsilon \int_ {\mathbf {R} _ {+}} e ^ {2 \pi (i x - \varepsilon) y} d y - 2 \pi \varepsilon \int_ {\mathbf {R} _ {-}} e ^ {2 \pi (i x + \varepsilon) y} d y \\ = \frac {\varepsilon}{\varepsilon - i x} - \frac {\varepsilon}{\varepsilon + i x} = \frac {2 i \varepsilon x}{x ^ {2} + \varepsilon^ {2}}, \end{array}
\]

由此 \(\overline{\mathcal{F}}(\eta_{\varepsilon}) = \sigma_{\varepsilon}\) 。于是结果由 \(\mathrm{L}^2(\mathbf{R})\) 中的傅里叶反演公式（II, p. 220 的定理 2 的推论）得出。

引理 7. --- 设 \(f \in \mathcal{D}(\mathbf{R})\) ，并设 \(\widetilde{f} \in \mathcal{D}(\mathbf{C})\) 是 \(f\) 的一个几乎解析延拓。对 \(\varepsilon > 0\) ，定义 \(f_{\varepsilon}\) 于 \(\mathbf{R}\) 上如下

\[
f _ {\varepsilon} (x) = - \frac {1}{2 i \pi} \int_ {\mathbf {R}} \biggl (\frac {\widetilde {f} (y + i \varepsilon)}{y - x + i \varepsilon} - \frac {\widetilde {f} (y - i \varepsilon)}{y - x - i \varepsilon} \biggr) d y.
\]

那么 \(f_{\varepsilon}\) 是连续且有界的，且 \(f_{\varepsilon}\) 收敛于 \(f\) （在 \(\mathcal{C}_b(\mathbf{R})\) 中，当 \(\varepsilon\) 趋于 0 时）。

\(f_{\varepsilon}\) 的连续性是 INT, IV, p. 144, § 4, 第 3 目, 推论 1 的推论，因为 \(\widetilde{f}\) 具有紧支撑。此外，若 \(r\) 使得 \(\widetilde{f}\) 的支撑包含在 \([-r, r] \times \mathbf{R}\) 中，我们有

\[
| f _ {\varepsilon} (x) | \leqslant 2 \| \widetilde {f} \| _ {\infty} \int_ {- r} ^ {r} \frac {1}{\sqrt {(x - y) ^ {2} + \varepsilon^ {2}}} d y \leqslant \frac {4 r}{\varepsilon} \| \widetilde {f} \| _ {\infty},
\]

因此 \(f_{\varepsilon}\) 是有界的。

一阶泰勒展开（FVR, I, p. 30, 命题 3）与公式 (13) 表明，存在 M ≥ 0 及一个函数 \(\varrho_{1}\) ，它取值于 \(R^{2}\) 中，使得

\[
\widetilde {f} (y + i \gamma) = f (y) + i \gamma f ^ {\prime} (y) + \gamma^ {2} \varrho_ {1} (y; \gamma), \text {   and   } | \varrho_ {1} (y; \gamma) | \leqslant M,
\]

对所有 \(y \in \mathbf{R}\) 及 \(\gamma \in \mathbf{R}\) 成立。由于映射 \(y \mapsto \widetilde{f}(y + i\gamma)\) 的支撑包含在 \([-r, r]\) 中（对所有 \(\gamma \in \mathbf{R}\) ），映射 \(y \mapsto \varrho_{1}(y; \gamma)\) 的支撑也包含在 \([-r, r]\) 中，对所有 \(\gamma \in \mathbf{R}\) 。

设 \(g_{\varepsilon}\) 是定义在 \(\mathbf{R}^2\) 上的函数

\[
g _ {\varepsilon} (x, y) = \frac {\widetilde {f} (y + i \varepsilon)}{y - x + i \varepsilon} - \frac {\widetilde {f} (y - i \varepsilon)}{y - x - i \varepsilon}.
\]

对每个 \((x,y)\in \mathbf{R}^2\) 及 \(\varepsilon >0\) ，我们得到

\[
\begin{array}{r l} g _ {\varepsilon} (x, y) = - \frac {2 i \varepsilon}{(x - y) ^ {2} + \varepsilon^ {2}} f (y) + \frac {2 i \varepsilon (y - x)}{(x - y) ^ {2} + \varepsilon^ {2}} f ^ {\prime} (y) \\ & \quad + \varepsilon^ {2} \bigg (\frac {\varrho_ {1} (y ; \varepsilon)}{y - x + i \varepsilon} - \frac {\varrho_ {1} (y ; - \varepsilon)}{y - x - i \varepsilon} \bigg). \end{array}
\]

设 \(x \in \mathbf{R}\) 且 \(\varepsilon > 0\) 。我们有

\[
\begin{array}{c} \left| \varepsilon^ {2} \int_ {\mathbf {R}} \left(\frac {\varrho_ {1} (y ; \varepsilon)}{y - x + i \varepsilon} - \frac {\varrho_ {1} (y ; - \varepsilon)}{y - x - i \varepsilon}\right) d y \right| \leqslant 2 \mathrm{M} \varepsilon^ {2} \int_ {- r} ^ {r} \frac {d y}{\sqrt {(x - y) ^ {2} + \varepsilon^ {2}}} \\ \leqslant 4 \mathrm{Mr} \varepsilon . \end{array}
\]

因此，对每个 \(x \in \mathbf{R}\) ，可得

\[
f _ {\varepsilon} (x) = - \frac {1}{2 i \pi} \int_ {\mathbf {R}} g _ {\varepsilon} (x, y) d y = (f * \delta_ {\varepsilon}) (x) - \frac {1}{2 i \pi} (f ^ {\prime} * \sigma_ {\varepsilon}) (x) + k _ {\varepsilon} (x)
\]

其中 \(\delta_{\varepsilon}\) 与 \(\sigma_{\varepsilon}\) 是 \(\mathbf{R}\) 上由下式定义的函数

\[
\delta_ {\varepsilon} (x) = \frac {1}{\pi} \frac {\varepsilon}{x ^ {2} + \varepsilon^ {2}}, \qquad \sigma_ {\varepsilon} (x) = \frac {2 i \varepsilon x}{x ^ {2} + \varepsilon^ {2}},
\]

且 \(\| k_{\varepsilon}\|_{\infty}\leqslant 2\mathrm{Mr}\varepsilon .\)

函数 \(\mathcal{F}(f') \in \mathcal{S}(\mathbf{R})\) 是可积的（IV, p. 219 的命题 18 及 IV, p. 213 的命题 13）。由引理 6，可得

\[
\begin{array}{r l r} & & {\int_ {\mathbf {R}} | \mathcal {F} (f ^ {\prime}) (y) \mathcal {F} (\sigma_ {\varepsilon}) (y) | d y = 2 \pi \varepsilon \int_ {\mathbf {R}} | \mathcal {F} (f ^ {\prime}) (y) | e ^ {- 2 \pi \varepsilon | y |} d y} \\ & & {\leqslant 2 \pi \varepsilon \int_ {\mathbf {R}} | \mathcal {F} (f ^ {\prime}) (y) | d y,} \end{array}
\]

因此 \(\mathcal{F}(f')\mathcal{F}(\sigma_{\varepsilon})\) 在 \(\mathrm{L}^{1}(\mathbf{R})\) 中收敛于 0，当 \(\varepsilon\) 趋于 0 时。由于 \(f'\) 与 \(\sigma_{\varepsilon}\) 属于 \(\mathrm{L}^{2}(\mathbf{R})\) ，II, p. 223 的命题 14 蕴含 \(f'*\sigma_{\varepsilon}=\overline{\mathcal{F}}(\mathcal{F}(f')\mathcal{F}(\sigma_{\varepsilon}))\) 在 \(\mathcal{C}_{b}(\mathbf{R})\) 中收敛于 0。

对每个 \(\varepsilon > 0\) ，正测度 \(\delta_{\varepsilon} \cdot dx\) 在 \(\mathbf{R}\) 上的总质量为 1（参见 FVR, III, p. 7）。测度 \(\delta_{\varepsilon} \cdot dx\) （对 \(\varepsilon > 0\) ）所成的集，与 \(\mathbf{R}_{+}^{*}\) 中 0 的邻域滤子在该集上诱导的滤子，满足 INT, VIII, p. 137, § 2, 第 7 目的引理 4 的假设。因此，测度 \(\delta_{\varepsilon} \cdot dx\) 在 \(\mathcal{M}^{1}(\mathbf{R})\) 中收敛于点测度 \(\varepsilon_{0}\) ，当 \(\varepsilon\) 趋于 0 时。由此可得 \(f * \delta_{\varepsilon} \to f\) 在 \(\mathcal{C}_{b}(\mathbf{R})\) 中成立（INT, VIII, p. 163, § 4, 第 4 目）。引理得证。

我们用 \(\mu\) 表示 \(\mathbf{C}\) 上的勒贝格测度。

定理 2（埃尔弗--舍斯特兰德）. --- 设 u 是 E 上的自伴部分算子。设 \(f \in \mathcal{D}(\mathbf{R})\) ，且 \(\widetilde{f} \in \mathcal{D}(\mathbf{C})\) 是 \(f\) 的一个几乎解析延拓。设 h 是由下式定义的从 \(\mathbf{C}\) 到 \(\mathcal{L}(\mathrm{E})\) 中的映射

\[
h (\lambda) = \frac {\partial \widetilde {f}}{\partial \overline {{z}}} (\lambda) \mathrm{R} (u, \lambda)
\]

若 \(\lambda\in\mathbf{C} - \mathbf{R}\) ；而 \(h(\lambda)=0\) ，若 \(\lambda\in\mathbf{R}\) 。那么 h 关于 \(\mu\) 在 \(\mathbf{C}\) 上可积，且

\[
f (u) = - \frac {1}{\pi} \int_ {\mathbf {C}} h (\lambda) d \mu (\lambda).
\]

映射 h 是可测的，且其支撑是紧的。由于 u 是自伴的，我们有 \(\|\mathrm{R}(u,\lambda)\| \leqslant |\mathcal{I}(\lambda)|^{-1}\) ，对所有 \(\lambda \in \mathbf{C} - \mathbf{R}\) （IV, p. 248 的命题 17）。因此，依据公式 (13)，映射 h 是有界的，从而它在 \(\mathbf{C}\) 上可积。

设 \(\varepsilon \in \mathbf{R}_{+}^{*}\) 。我们用 \(\mathrm{F}_{\varepsilon}^{+}\) （相应地，\(\mathrm{F}_{\varepsilon}^{-}\) ）表示 \(\mathbf{C}\) 中由这样的 \(\lambda \in \mathbf{C}\) 所成的闭集：\(\mathcal{I}(\lambda) \geqslant \varepsilon\) （相应地，\(\mathcal{I}(\lambda) \leqslant -\varepsilon\) ）。我们有

\[
\int_ {\mathbf {C}} h (\lambda) d \mu (\lambda) = \lim _ {\varepsilon \to 0} \Bigl (\int_ {\mathrm{F} _ {\varepsilon} ^ {+}} h (\lambda) d \mu (\lambda) + \int_ {\mathrm{F} _ {\varepsilon} ^ {-}} h (\lambda) d \mu (\lambda) \Bigr).
\]

设 r > 0 使得 h 的支撑包含在 \(\mathbf{C} = [-r, r]^{2}\) 中。用 \(R_{\varepsilon}^{+}\) 表示 \([-2r, 2r] \times [\varepsilon, 2r]\) 这一 \(\mathbf{C}\) 中的长方体。它是 \(\mathbf{C}\) 的局部多面体子集（VAR, R2, 11.3, p. 48）。它满足下列条件：

\begin{enumerate}
\def\labelenumi{(\roman{enumi})}
\item
  我们有 \(R_{\varepsilon}^{+} \subset 2C \cap F_{\varepsilon}^{+}\) ；
\item
  集合 \(R_{\varepsilon}^{+}\) 包含 \(F_{\varepsilon}^{+}\) 与 h 的支撑的交；
\item
  正则边界 \(\partial R_{\varepsilon}^{+}\) （VAR, R2, 11.3.2, p. 49）包含线段 \(S_{\varepsilon} = [-r, r] + i\varepsilon \subset C\) ；
\item
  有 \(h(\lambda) = 0\) ，只要 \(\lambda \in \partial \mathrm{R}_{\varepsilon}^{+} - \mathrm{S}_{\varepsilon}\) 。
\end{enumerate}

用 \(d\lambda\) （相应地，\(d\overline{\lambda}\) ）表示 \(\mathbf{C}\) 上的、\(\mathbf{C}\) 的恒等映射（相应地，复共轭）的微分 1-形式。用 \(g\) 表示 \(\mathbf{C} - \mathbf{R}\) 上取值于 \(\mathcal{L}(\mathrm{E})\) 的、使得 \(g(\lambda) = \widetilde{f}(\lambda)\mathrm{R}(u,\lambda)\) （对 \(\lambda \in \mathbf{C} - \mathbf{R}\) ）的函数。设 \(\omega = g d\lambda\) ；它是 \(\mathbf{C} - \mathbf{R}\) 上具有紧支撑且取值于 \(\mathcal{L}(\mathrm{E})\) 的 1 次微分形式。由于 \(u\) 的预解式是全纯的（IV, p. 246 的命题 14），我们有

\[
d \omega = \left(\frac {\partial \widetilde {f}}{\partial \overline {{z}}} (\lambda) \mathrm{R} (u, \lambda) + \widetilde {f} (\lambda) \frac {\partial}{\partial \overline {{z}}} \mathrm{R} (u, \lambda)\right) d \overline {{\lambda}} \wedge d \lambda = - h (\lambda) d \lambda \wedge d \overline {{\lambda}}.
\]

与 \(d\omega\) 相伴的向量测度（VAR, R2, 10.4.3, p. 43）是关于勒贝格测度以 -2ih 为密度的测度。对局部多面体集 \(R_{\varepsilon}^{+}\) 及微分形式 \(\omega\) 应用斯托克斯公式（VAR, R2, 11.3.4, p. 49），我们便得到

\[
\frac {i}{2} \int_ {\mathrm{R} _ {\varepsilon} ^ {+}} d \omega = \frac {i}{2} \int_ {\partial \mathrm{R} _ {\varepsilon} ^ {+}} \omega = \frac {i}{2} \int_ {\mathrm{S} _ {\varepsilon}} \omega = \frac {i}{2} \int_ {- r} ^ {r} \widetilde {f} (y + i \varepsilon) \mathrm{R} (u, y + i \varepsilon) d y,
\]

由此

\[
\int_ {\mathrm{F} _ {\varepsilon} ^ {+}} h (\lambda) d \mu (\lambda) = - \frac {i}{2} \int_ {\mathbf {R}} \widetilde {f} (y + i \varepsilon) \mathrm{R} (u, y + i \varepsilon) d y.
\]

对 \(F_{\varepsilon}^{-}\) 作同样的推理，我们得到

\[
\int_ {\mathrm{F} _ {\varepsilon} ^ {-}} h (\lambda) d \mu (\lambda) = - \frac {i}{2} \int_ {\mathbf {R}} \widetilde {f} (y - i \varepsilon) \mathrm{R} (u, y - i \varepsilon) d y,
\]

并且我们得出，\(h\) 在 \(\mathbf{C}\) 上的积分是当 \(\varepsilon\to0\) 时下列量的极限

\[
v _ {\varepsilon} = \frac {i}{2} \int_ {\mathbf {R}} \Bigl (\widetilde {f} (y + i \varepsilon) \mathrm{R} (u, y + i \varepsilon) - \widetilde {f} (y - i \varepsilon) \mathrm{R} (u, y - i \varepsilon) \Bigr) d y.
\]

由命题 7，我们有 \(v_{\varepsilon} = \pi f_{\varepsilon}(u)\) ，其中 \(f_{\varepsilon}\) 是定义在 \(\mathbf{R}\) 上的函数

\[
f _ {\varepsilon} (x) = - \frac {1}{2 i \pi} \int_ {\mathbf {R}} \biggl (\frac {\widetilde {f} (y + i \varepsilon)}{y - x + i \varepsilon} - \frac {\widetilde {f} (y - i \varepsilon)}{y - x - i \varepsilon} \biggr) d y.
\]

由于 \(f_{\varepsilon} \to f\) 当 \(\varepsilon \to 0\) 时在 \(\mathcal{C}_{b}(\mathbf{R})\) 中成立（引理 7），自同态 \(v_{\varepsilon} = \pi f_{\varepsilon}(u)\) 收敛于 \(\pi f(u)\) ，在 \(\mathcal{L}(\mathrm{E})\) 中（IV, p. 275 的命题 5）。定理得证。

